\documentclass[11pt,a4paper]{ltjsarticle}

% 日本語版．英語正本：02-v5-mcfg-regular-skeletons/main.tex
% LuaLaTeX でコンパイルする：lualatex main_JP.tex（2回）
% ラベルは英語版と同一に保つ（tools/check_labels.py で検査）．

\usepackage{amsmath,amssymb,amsthm,mathtools}
\numberwithin{equation}{section}
\usepackage{enumitem}
\usepackage{geometry}
\geometry{margin=1in}
\usepackage{booktabs}
\usepackage{cite}
\usepackage{hyperref}
\usepackage{url}
\hypersetup{hidelinks}
\setlength{\emergencystretch}{3em}

\theoremstyle{plain}
\newtheorem{theorem}{定理}[section]
\newtheorem{lemma}[theorem]{補題}
\newtheorem{proposition}[theorem]{命題}
\newtheorem{corollary}[theorem]{系}
\newtheorem{fact}[theorem]{事実}
\theoremstyle{definition}
\newtheorem{definition}[theorem]{定義}
\newtheorem{example}[theorem]{例}
\theoremstyle{remark}
\newtheorem{remark}[theorem]{注意}
\newtheorem{openproblem}[theorem]{未解決問題}
\theoremstyle{plain}
\newtheorem{mainthm}{定理}
\renewcommand{\themainthm}{\Alph{mainthm}}

\newcommand{\Cmcf}[2]{\mathcal{C}^{\mathrm{mcf}}_{#1,#2}}
\newcommand{\D}{\mathcal{D}}
\newcommand{\E}{\mathcal{E}}
\newcommand{\tuple}[1]{\vec{#1}}
\newcommand{\blank}{\square}
\newcommand{\eps}{\varepsilon}
\newcommand{\alphset}{\operatorname{alph}}
\newcommand{\Imp}{\operatorname{Imp}}
\newcommand{\fan}{\operatorname{fan}}
\newcommand{\LG}{\mathbf{LG}}
\newcommand{\LI}{\mathbf{LI}}
\newcommand{\SL}{\mathbf{SL}}
\newcommand{\Gp}{\mathbf{G}}
\newcommand{\Dd}{\mathbf{D}}
\newcommand{\ECom}{\mathbf{ECom}}
\newcommand{\LECom}{\mathbf{L}\mathbf{ECom}}
\newcommand{\LT}{\mathbf{LT}}
\newcommand{\LTT}{\mathbf{LTT}}
\newcommand{\ACom}{\mathbf{ACom}}
\newcommand{\Com}{\mathbf{Com}}
\newcommand{\ttyp}{t_{\mathrm{typ}}}
\newcommand{\bth}{\bar\theta}

\title{正則な骨格を持つ多重文脈自由言語の多項式特性データ\\
——局所群の外の型付けへ}
\author{栗山貴之}
\date{草稿，2026年11月23日}

\begin{document}
\maketitle

\begin{abstract}
固定した有限半群型付け $s$ のもとで，正データから多重文脈自由言語を分布学習するとき，特性データの総長は，標的の表示 $G$ の\emph{代表長} $\tau(G,s)$ で制御される．これは，型付き非終端と同じ型・同じ分布を持つ最短の組の長さである．型付けの像が局所群（擬多様体 $\LG$）なら，$\tau$ は全 fan-out で $G$ の大きさと厚さの多項式であり，fan-out 1 では $\LG$ が型付けについてのちょうどの境界である．$\LG$ の外については，fan-out 1 を越えるとほとんど何も知られていなかった．

本論文はこの問いを標的の側から扱う．標的を，正則な\emph{骨格} $R$ と，長さが有界な因子についての任意の加法的な計数条件とで記述する．$R$ の統語半群のすべての局所モノイドで冪等元が可換なら，すなわち統語半群が $\LECom$ に入るなら，内容・固定幅の窓・$R$ の統語射を細分する任意の型付け（標的がその型付けで代入可能であることを仮定する）について，固定した fan-out と rank のもとで，任意の reduced good な表示の非終端数と厚さの多項式で代表長が抑えられることを示す．これらの型付けは一般に $\LG$ の外にある．証明では，遷移半群が局所群であるオートマトンの状態で各文字を装飾し，装飾後のアルファベット上の局所計数言語に標的を移し，局所計数言語についての鳩の巣と Euler 再配線の議論を適用する．装飾からどの入口でも統語値全体を読み取れる骨格は，統語半群が $\SL*\LG$ に入るものにちょうど一致し，$\SL*\LG=\LECom$ である．さらに，この方法がちょうどそこで止まることも示す．$\LECom$ の外の骨格では，局所群による装飾のあとどんな局所計数条件を課しても所属は決まらず，最小の例を三つ与える．型付けの側の境界（$\LG$：局所モノイドが群）と骨格の側の境界（$\LECom$：局所モノイドの冪等元が可換）は，こうしてどちらも局所モノイドについての条件になる．
\end{abstract}

\tableofcontents

%======================================================================
\section{序論}
\label{sec:intro}

\paragraph{背景．}
吉仲 \cite{Yoshinaka2011} は，多次元の代入可能性を持つ多重文脈自由言語（MCFL）\cite{SekiEtAl1991} が，正データから極限同定可能であり，更新は多項式時間，特性標本の要素数は多項式であることを示した．先行研究 \cite{KuriyamaTCS1,KuriyamaMCFGpreprint,KuriyamaCompendium} では，学習器を固定した有限半群型付け $s:\Sigma^+\to\mathsf S$ で偏らせる．二つの組は，成分の $s$-型が等しいときだけ比べられる．$(f,s)$-組代入可能な $f$-MCFL のクラス $\Cmcf{f}{s}$ は，$s$ が認識するすべての正則言語を含み，$s$ とともに大きくなる．

この一般性の代償は，特性データの総長で測られる．標的の reduced good な表示 $G$ には\emph{代表長} $\tau(G,s)$ がある．実現された型付き非終端 $X$ ごとに，$X$ の生成組と同値な（同じ型・同じ分布の）最短の組の長さを取り，$X$ について最大を取ったものである．総長 $O(N_{\mathrm{typ}}(1+\bth_G+r\tau(G,s)))$ の特性データが存在する（事実~\ref{fact:char-data}）ので，$\tau$ を $G$ の大きさと厚さの多項式で抑えれば，特性データは多項式になる．

\emph{型付け}の側の状況は次のとおりである．$s(\Sigma^+)$ が局所群（どの局所モノイド $e\mathsf Se$ も群，擬多様体 $\LG$）なら，型付き生成組そのものが全 fan-out で短い \cite{KuriyamaCompendium}．fan-out 1 では $\LG$ がちょうど擬多様体の境界である．$\LG$ の外の半群への型付けには，特性データが指数になる標的がある \cite{KuriyamaTCS1,KuriyamaTypingClass}．fan-out 2 以上で $\LG$ がなお境界かどうかは未解決である（未解決問題~\ref{op:fanout-two}）．

\paragraph{本論文：標的の側．}
型付けを制限する代わりに，標的を制限する．標的を，正則な\emph{骨格} $R$（その統語射 $\rho_R$ を通して）と，固定長 $\ell$ の因子の個数 $N_\ell(w)$ についての任意の加法的な条件とで記述する．型付けには，内容・固定幅の窓・$\rho_R$ が見えることを要求する．これは一般に $\LG$ に入らない．

\begin{mainthm}[主定理；定理~\ref{thm:main}]
\label{thm:A}
$R\subseteq\Sigma^+$ を正則とし，その統語半群が $\LECom$ に入る，すなわち $\Sigma^+/{\sim_R}$ のどの局所モノイドでも任意の二つの冪等元が可換であるとする．$\ell,f,r\ge1$ を固定する．正の語の $L$ への所属が $(\rho_R(w),N_\ell(w))$ だけで決まり，$s\preceq(c_\Sigma,b_\ell,\rho_R)$，$L\in\Cmcf{f}{s}$ とする．このとき，$L\cap\Sigma^+$ の任意の reduced good，fan-out $\le f$，rank $\le r$ の表示 $G$ について
\[
  \tau(G,s)\le\mathrm{poly}(v,\bth_G)
\]
が成り立つ．ここで $v$ は非終端の数で，次数は $R,\Sigma,\ell,f$ だけで決まる．したがって $L$ は総長が多項式の特性データを持つ．
\end{mainthm}

ここで $c_\Sigma$ は内容の型付け，$b_\ell$ は最初と最後の $\ell-1$ 文字を記録する窓の型付けである．定理~\ref{thm:A} の型付けは，$|\Sigma|\ge2$ ならすでに $\LG$ の外にある（補題~\ref{lem:window-not-lg}）．$R$ の仮定は，たとえば，すべての局所閾値テスト可能言語，すべての群言語，そして極大な $\mathtt a$-run の長さがすべて偶数である語の言語で満たされる．最後のものの統語半群は対称逆モノイド $I_2\notin\LG$ である（例~\ref{ex:even-runs}）．

\paragraph{方法．}
証明は三層からなる．
\begin{enumerate}[label=(\arabic*),leftmargin=*]
\item \emph{局所計数の標的}（定理~\ref{thm:L}）．所属が窓の型と因子の個数だけで決まるなら，de Bruijn グラフの辺の頻度についての鳩の巣，深さ優先の走査に沿った不純な窓の色，凍結ブロック付きの Euler 再配線から，短い同値な代表元が得られる．
\item \emph{装飾}（定理~\ref{thm:B}）．固定したオートマトンの状態で各文字を装飾すると，骨格は大きなアルファベット上の局所計数言語になる．装飾後の表示はまた reduced good で，オートマトンの遷移半群が局所群なら厚さは多項式のままである．元の分布は射影で回復する．
\item \emph{骨格の代数}（定理~\ref{thm:C}）．どの入口からでも統語値全体を読み取れる局所群の装飾を骨格が持つのは，統語半群が $\SL*\LG$ に入るときに限り，しかも $\SL*\LG=\LECom$ である．
\end{enumerate}

\begin{mainthm}[方法の限界；定理~\ref{thm:limits}]
\label{thm:D}
正則な骨格 $R$ の統語半群のある局所モノイドに，非可換な二つの冪等元があるとする．このとき，遷移半群が局所群である任意の固定オートマトン，任意の幅 $m$，装飾語の端の窓と因子の個数についての任意の条件について，それらのデータでは $R$ への所属は決まらない．最小の三例は，最後の $\mathtt a$-run の長さが偶数であるもの，中立な文字を持つ散在パターン $\mathtt a\cdots\mathtt b$，恒等の文字を持つ二状態の flip-flop である．
\end{mainthm}

\paragraph{局所モノイドについての二つの境界．}
\begin{center}
\begin{tabular}{lll}
\toprule
側 & 境界 & どの局所モノイド $e\mathsf Se$ についての条件\\
\midrule
型付け（fan-out 1；\cite{KuriyamaTCS1,KuriyamaTypingClass}） & $\LG$ & 群である\\
標的の骨格（本論文） & $\LECom$ & 冪等元が可換\\
\bottomrule
\end{tabular}
\end{center}
$\LG\subsetneq\LECom$ なので，後者は前者から従わない．どちらも，型付けのすべての標的について問う未解決問題~\ref{op:fanout-two} には答えない．関係は第~\ref{sec:limits}節で述べる．

\paragraph{構成．}
第~\ref{sec:prelim}節で記法を定める．第~\ref{sec:L}節で局所計数の標的についての定理，第~\ref{sec:decoration}節で移送定理，第~\ref{sec:algebra}節で代数的な特徴づけ，第~\ref{sec:main}節で主定理と系を示す．第~\ref{sec:limits}節で方法の限界を論じ，第~\ref{sec:conclusion}節で結ぶ．

%======================================================================
\section{準備}
\label{sec:prelim}

\subsection{語，窓，型付け}

$\Sigma$ は有限アルファベットで $k=|\Sigma|$，$\Sigma^+$ は正の語の自由半群である．\emph{型付け}とは有限半群への全射準同型 $s:\Sigma^+\to\mathsf S$ のことで，$s^{(d)}$ は $d$ 組に成分ごとに $s$ を施す．型付け $s,h$ について，$h(x)=h(y)$ から $s(x)=s(y)$ が従うとき $s\preceq h$ と書く．型付けの組は $(s,h)(x)=(s(x),h(x))$ で合わせる．

\emph{内容}の型付けは $c_\Sigma(w)=\alphset(w)$ である．$j\ge1$ について $q=j-1$ とおく．\emph{窓}の型付け $b_j$ は，$|w|<\max(1,2q)$ なら語 $w$ そのものに短い印を付けたものを，そうでなければ長さ $q$ の接頭辞と接尾辞の対に印を付けたものを返す．連結の型は各片の型で決まるので，$b_j$ は（有限半群への）型付けである．$N_j(w)\in\mathbb N^{\Sigma^j}$ は，$w$ の長さ $j$ の因子を重なりを許してすべて数える．

長さ $q$ 以上の語は，頂点集合 $\Sigma^q$，辺集合 $\Sigma^j$ の \emph{de Bruijn グラフ}の中の，長さ $q$ の接頭辞から接尾辞への路に対応する．語は最初の $q$ 文字に続けて各辺の最後の文字を並べたものである．$N_j$ は辺を数える．

\subsection{多重文脈自由文法}

多重文脈自由文法（MCFG）\cite{SekiEtAl1991} を \cite{Yoshinaka2011} の記法で用いる．規則 $A\to f[B_1,\ldots,B_n]$ は，$A$ の各出力成分を，終端記号と子の成分の変数を連結して作る．規則が \emph{good} であるとは，テンプレートの成分が空でなく，各変数がちょうど一回現れ，各子の変数が連結した出力の中で添字順に現れ，一つの出力成分の中に同じ子の連続する二変数が隣接しないことをいう．開始記号の fan-out は 1 である．すべての非終端が到達可能かつ生成的なとき文法は \emph{reduced} である．非終端の数を $v=|V|$，子の数の上限（rank）を $r$，fan-out の上限を $f$ と書く．good な文法では，生成される組の成分はすべて正である．

\emph{厚さ} $\bth_G\ge1$ は，1 と，規則ごとの「その規則を使う成功文の最短長」の最大との大きい方である．各規則が書く終端記号は $\bth_G$ 個以下で，どの非終端も総長 $\bth_G$ 以下の生成組を持つ．

\subsection{文脈，代入可能性，標的のクラス}

arity $d$ の\emph{許容文脈}とは，$v_0,v_d\in\Sigma^*$，$v_1,\ldots,v_{d-1}\in\Sigma^+$ による語 $E=v_0\blank v_1\cdots v_{d-1}\blank v_d$ であり，その集合を $\E_d$ と書く．$\tuple x\in(\Sigma^+)^d$ の\emph{分布}は $\D_L(\tuple x)=\{E\in\E_d:E[\tuple x]\in L\}$ である．

\begin{definition}
\label{def:class}
$1\le d\le f$ と $\tuple x,\tuple y\in(\Sigma^+)^d$ について，$s^{(d)}(\tuple x)=s^{(d)}(\tuple y)$ かつ $\D_L(\tuple x)\cap\D_L(\tuple y)\ne\emptyset$ ならつねに $\D_L(\tuple x)=\D_L(\tuple y)$ となるとき，$L\subseteq\Sigma^*$ は \emph{$(f,s)$-組代入可能}であるという．$(f,s)$-組代入可能な $f$-MCFL のクラスを $\Cmcf{f}{s}$ と書く．$s^{(d)}(\tuple x)=s^{(d)}(\tuple y)$ かつ $\D_L(\tuple x)=\D_L(\tuple y)$ のとき $\tuple x\equiv_{L,s}\tuple y$ と書く．
\end{definition}

$s\preceq s'$ なら $\Cmcf{f}{s}\subseteq\Cmcf{f}{s'}$ である．このクラスは，可換群の剰余類の逆像との交わりで閉じている \cite[命題 \textit{abelian-coset-filter}]{KuriyamaCompendium}．すなわち，$L$ が $(f,s)$-組代入可能，$g:\Sigma^*\to\mathsf A$ が可換群への準同型，$\mathsf H\le\mathsf A$，$\gamma\in\mathsf A$ なら，$L\cap g^{-1}(\gamma+\mathsf H)$ もまた $(f,s)$-組代入可能である．

\paragraph{型付き非終端と代表長．}
\emph{実現された型付き非終端} $X=(A,\mathbf t)$ とは，非終端 $A$ と，成功する導出の中に現れる出力型のベクトル $\mathbf t$ の組である．$A$ の到達可能な出現は，その生成組すべてに共通の受理文脈を持つ（兄弟は正で，good な規則は穴の順序と分離を保つ）．したがって $L\in\Cmcf{f}{s}$ なら，$X$ の生成組はすべて一つの $\equiv_{L,s}$-類 $\operatorname{Cl}(X)$ に入る．
\[
  \tau_X=\min\{\|\tuple y\|_1:\tuple y\in\operatorname{Cl}(X)\},\qquad
  \tau(G,s)=\max_X\tau_X
\]
とおく．最小を与える組は $A$ から生成される必要はない．\emph{型付き厚さ} $\ttyp(G,s)$ は，$\tuple y$ を $X$ の生成組に限って同様に定める．したがって $\tau\le\ttyp$ である．

\begin{fact}[{\cite[定理 \textit{data-from-typed-thickness}，系 \textit{distribution-representatives}]{KuriyamaCompendium}}]
\label{fact:char-data}
$L\in\Cmcf{f}{s}$ とし，$G$ を $L\cap\Sigma^+$ の reduced good，fan-out $\le f$，rank $\le r$ の表示とする．このとき $L$ は，\cite{KuriyamaMCFGpreprint,KuriyamaCompendium} の学習器について，$\|K\|_+\le N_{\mathrm{typ}}\bigl(1+2\bth_G+r\,\tau(G,s)\bigr)+\mathbf 1_{\{\eps\in L\}}$ の特性標本 $K$ を持つ．ここで $N_{\mathrm{typ}}\le|\mathsf S^\sharp|^{fr}\|G\|$ である．
\end{fact}

\subsection{擬多様体}

\emph{擬多様体}とは，割り算と有限直積で閉じた有限半群（またはモノイド）のクラスである．次を用いる：$\Gp$（群），$\SL$（半束），$\Com$ と $\ACom$（可換モノイド，可換で非周期的なモノイド），$\Dd$（冪等元が右零元である半群），$\LI$（局所自明），$\LT$ と $\LTT$（局所テスト可能言語と局所閾値テスト可能言語の統語半群），$\ECom$（冪等元が可換なモノイド）．モノイドの擬多様体 $\mathbf V$ について，$\mathbf L\mathbf V$ はすべての局所モノイド $e\mathsf Se$（$e=e^2$）が $\mathbf V$ に入る有限半群のクラスである．したがって $\LG=\mathbf L\Gp$ であり，$\LECom$ はすべての局所モノイドで冪等元が可換な有限半群，すなわち擬恒等式
\[
  (exe)^\omega(eye)^\omega=(eye)^\omega(exe)^\omega\qquad(e=z^\omega)
\]
を満たす有限半群からなる．

モノイド $U$ と，集合 $Q$ に右から作用する半群 $T$ について，\emph{花輪積} $U\wr_QT$ は $F:Q\to U$，$t\in T$ の対 $(F,t)$ からなり，積は $(F,t)(H,t')=(q\mapsto F(q)H(qt),tt')$ である．擬多様体について，$\mathbf V*\mathbf W$ は，$U\in\mathbf V$，$T\in\mathbf W$ によるこうした花輪積を割り切る半群で生成される．右の因子が先に入力を読む．この演算は結合的である \cite{Tilson1987}．有限半群 $T$ が $\LG$ に入るのは，$T$ が冪等元を持たないか，極小イデアルが完全単純で $T$ のすべての冪等元を含むときに限る \cite[命題2.1]{CostaNogueiraTeixeira2021}．

\begin{lemma}
\label{lem:window-not-lg}
$k\ge2$，$\ell\ge1$ なら，$(c_\Sigma,b_\ell)$ の像は $\LG$ に入らない．
\end{lemma}

\begin{proof}
$\mathtt a\ne\mathtt b$ を文字とし，$e$ を $\mathtt a^{2\ell}$ の型とする．$e$ と $e\,s(\mathtt b)\,e$ は局所モノイド $e\mathsf Se$ の相異なる冪等元である．端の窓は同じで内容が違う．群の冪等元は一つしかない．
\end{proof}

%======================================================================
\section{局所計数の標的}
\label{sec:L}

この節では $\ell\ge1$，$q=\ell-1$，$V=k^q$，$K=k^\ell$，$b=b_\ell$，$N=N_\ell$，$c=c_\Sigma$，$T=|b(\Sigma^+)|$ とする．

\begin{definition}
正の語の所属が $(b_\ell(w),N_\ell(w))$ だけで決まるとき，$L$ は \emph{$\ell$-局所計数}であるという．
\end{definition}

位数 $\ell$ の厳密局所テスト可能・局所テスト可能・局所閾値テスト可能な言語に，文字の個数についての任意の条件を交わらせたものは $\ell$-局所計数である．文字の個数は，最初の頂点と辺の最後の文字から読み取れる．

\begin{theorem}[局所計数の標的]
\label{thm:L}
$L\in\Cmcf{f}{s}$ を $\ell$-局所計数とし，$s\preceq(c,b)$ とする．
\[
 \lambda=K(2V-1),\quad B=2q\lambda,\quad c_1=1,\quad c_{i+1}=Kc_i+f+B+2
 \ (1\le i\le K+1),\quad c_*=c_{K+2}
\]
とおき，$L\cap\Sigma^+$ の reduced good，fan-out $\le f$，rank $\le r$，非終端数 $v$，$\theta=\bth_G$ の表示 $G$ について
\[
 H=v\bigl((2^k-1)T+1\bigr)^f,\quad N_{\rm sk}=\ell c_*KfH,\quad
 C_L=T^f\bigl(1+rH(\theta+f\ell)\bigr)^{Kf}
\]
とおく．このとき
\[
 \tau(G,s)\le f\ell+\ell N_{\rm sk}\theta+2\ell frN_{\rm sk}
 +\ell rN_{\rm sk}C_LH\theta
 =O_{k,\ell,f,r}\bigl(v^{Kf+2}\theta^{Kf+1}\bigr).
\]
\end{theorem}

\begin{proof}
\emph{文脈の不変量．}
文脈 $E[\tuple x]$ の境界をまたぐ窓は，最初に触れる片の接尾辞，最後に触れる片の接頭辞，その間のすべての空でない片からなる．間の片は長さ $q-1$ 以下なので，印で決まる．したがって，成分の $(b,c)$-型が等しく，内部の $N$ の和が等しければ，任意の許容文脈 $E$ について $E[\tuple x]$ の $(b,N)$ は等しく，分布も等しい．固定した任意の片についても同じ恒等式が成り立つ．

\emph{木と位置．}
実現された型付き非終端 $X$ を固定し，出力ラベルが $X$ の導出のうち節点数が最小のものを選び，その組を $\tuple x$ とする．記号と出力の $(c,b)$-型が等しい祖先と真の子孫があれば縮められるので，節点で数えた高さは $H$ 以下である．位置はつねにこの元の組について言う．部分木の各成分は一つの根成分の中の連続した区間に写り，比較不能な部分木の像は交わらない．

長さ $\ell$ 以上の根成分 $x_i$ ごとに，辺の頻度の $K+1$ 個の区間のうち空のもの $[c_g,c_{g+1})$ を一つ選ぶ．頻度 $c_{g+1}$ 以上の辺を\emph{頻出}，それ以外を\emph{稀}と呼ぶ．頻出の辺があれば，稀な出現は $R=(K-1)(c_g-1)$ 個以下である．頻出の辺の弱連結成分 $J$ はどれも強連結である．戻れない辺があれば，その終点から頻出の辺で到達できる集合はその始点を含まず，その集合に $c_{g+1}$ 回以上入るには $R$ 個を超える稀な出口が要るからである．頻出の辺の極大区間はどれも一つの $J$ に入り，各 $J$ の区間は $R+1$ 個以下である．

$x_i$ の各辺 $e$ について，最初の $\min(c_*,N(x_i)(e))$ 個の出現のすべての文字に印を付ける．印付きの文字を書く節点とその祖先を\emph{骨格}と呼ぶ．骨格の節点は $N_{\rm sk}$ 個以下である．骨格が空なら窓はなく $\|\tuple x\|_1\le fq$ である．そうでなければ，骨格の外の極大部分木は $rN_{\rm sk}$ 個以下である．それらの成分の内部の窓はどれも印に触れず，すべて一つの頻出区間に入る．短い印を持つ根成分は，すべての窓が選ばれていて，変更しない．

\emph{切り口と色．}
骨格の外の極大部分木 $\eta$（成分 $z_\kappa$）を固定し，その葉を深さ優先で並べる．最初の $j$ 枚の葉のあとで，葉がすべてそれらに含まれる極大部分木の集合を $I_j$ とする．$|I_j|\le rH$ で，その元の終端位置は増大列 $U_j$ をなし，一歩で増える位置は $H\theta$ 個以下である．$I_j$ の各元を，それ自身の記号から生成される長さ $\theta$ 以下の組の導出で置き換える．新しい成分を $z^j_\kappa$，得られる文を $W_j$ と書く．$W_j$ は生成される．挿入した文字すべてに印を付け，$j$ に\emph{色} $(b(z^j_\kappa),\Imp_\kappa(j))_\kappa$ を与える．$\Imp_\kappa(j)$ は，$z^j_\kappa$ の内部の窓のうち挿入文字に触れるもののベクトルである．置き換える片はすべて正なので，挿入文字を含まない窓は一つの残った区間の中にあり，位置によって $U_j$ を避ける元の内部の窓と一対一に対応する．色は $C_L$ 種類以下である．$j<j'$ の色が等しければ
\[
 Z_\kappa=N(z^j_\kappa)-N(z^{j'}_\kappa)\ge0
\]
は，$(j,j']$ の間に初めて触れられた元の内部の窓を数える．端の印が等しいので $Z_\kappa$ は一つの $J$ に台を持つ循環流であり，文脈の不変量から $b(W_j)=b(W_{j'})$，$N(W_j)-N(W_{j'})=\sum_\kappa Z_\kappa$ である．訪れた各添字から，同じ色の最後の添字へ跳び，その間の初接触を削除し，次の添字への増分一つだけを残して続ける．跳躍の窓は互いに交わらない．

\emph{Euler 再配線．}
強連結な $J$ の中で，$m$ 本の指定された路が端点 $p_l,q_l$ を持ち，合計の目標個数 $M$ が $\partial M=\sum_l(\mathbf e_{q_l}-\mathbf e_{p_l})$ を満たすとする．指定された主路の中の頂点 $u$ に閉歩道 $g_0$ を露出させるには，$J$ 上で $M(e)\ge m+1+N(g_0)(e)$ であれば足りる．主路以外の枠と，主路の始点から $u$ までに単純路を選び，それらと $g_0$ を引く．残りはすべての辺で正で，$u$ から主路の終点へ釣り合っているので Euler 路を持つ．指定された路の間の凍結された辺は，端点が保たれるので本物の窓のまま残る．

\emph{一回の削除．}
各根成分について，元の $(b,c)$ と，それまでの跳躍を引いた現在のベクトル $M_i$ を保ち，実現する組が存在してそのすべてが $\tuple x$ と同値である，という不変量を保つ．零でない $Z_\kappa$ を根成分と頻出成分ごとにまとめて溜まり $Z_{i,J}$ とする．活動中の溜まりは $f$ 個以下である．$W_{j'}$ では，$\eta$ の区間の外にある $x_i$ の内部の頻出の窓だけを動かしてよいとし，ほかの辺はすべて凍結する．柔軟な連は $R+1+f$ 本以下で，頻出の辺はどれも $c_{g+1}$ 個以上の柔軟な出現を持つ（選ばれた出現はすべて印付きなので $\eta$ の外にある）．活動中の各溜まりについて，柔軟な窓を含む最初の元の頻出区間を錨とする．

$u\in J$ と，$u$ から出て $J$ のすべての辺を通る長さ $\lambda$ 以下の巡回 $D$ を固定する．$q>0$ なら $D_0=D^q$，$h=D_0D_0$，$t=D_0\operatorname{Euler}_u(N(D_0)+Z_{i,J})$ とし，$q=0$ なら $h$ は $J$ の各ループを一回ずつ通り，$t$ は $N(h)+Z_{i,J}$ を通るとする．出力される語 $h,t$ は正で，内容も窓の型も等しく，$N(t)-N(h)=Z_{i,J}$ である．

代表元の側では，活動中の各溜まりを再配線して錨に $t$ を露出させ，$\tuple y$ を得る．$t$ を $h$ に替えると，個数がちょうど跳躍の分だけ減った $\tuple y'$ が得られる．文の側では，錨の柔軟な連に $h$ を露出させ，$W_{j'}$ の個数を持つ $F[h]$ と，$W_j$ の個数を持つ $F[t]$ を得る．個数の条件は，辺ごとに $m\le R+1+f$ と
\[
 c_{g+1}-(R+1+f+B+1)=c_g+K-1\qquad(q>0)
\]
から従う．$N(h)(e)\le B$ であり，代表元の側では $Z_{i,J}(e)$ が $t$ の分を払う．$q=0$ なら余裕は 1 以上である．$F[h]$ と $F[t]$ はともに $L$ に入る．一つの根成分の中の穴は凍結された稀な辺で，違う成分の穴は正の文脈片で隔てられるので，$F$ と $F_E$ は任意の許容文脈 $E$ について許容である．比べる組は $t$ と $h$，共有文脈は $F$，試験文脈は $F_E$ で，四つの語は $F[t],F[h],F_E[t]=E[\tuple y],F_E[h]=E[\tuple y']$ である．代入可能性から，すべての $E$ について $E[\tuple y]\in L\iff E[\tuple y']\in L$ となり，文脈の不変量でこれを新しいデータのすべての実現に広げる．

\emph{長さ．}
一つの自由成分の内部にない残った窓は，骨格が書く終端記号に触れる（始点は $\ell N_{\rm sk}\theta$ 個以下）か，自由成分の端をまたぐ（$2\ell frN_{\rm sk}$ 個以下）．ほかの残った窓は，その初接触の増分に割り当てる．自由部分木ごとに残す増分は $C_L$ 個以下で，それぞれが触れる窓は $\ell H\theta$ 個以下である．各成分の長さは $q$ と窓の数の和以下である．
\end{proof}

\begin{corollary}
\label{cor:SL-coset}
$R$ を厳密 $\ell$-局所テスト可能，$g:\Sigma^*\to\mathsf A$ を可換群への準同型，$\mathsf H\le\mathsf A$，$\gamma\in\mathsf A$ とする．$L=R\cap g^{-1}(\gamma+\mathsf H)$ が $f$-MCFL なら，$L\in\Cmcf{f}{(c,b_\ell)}$ であり，定理~\ref{thm:L} が適用できる．
\end{corollary}

\begin{proof}
arity $d\le f$ の，型が等しい正の組 $\tuple x,\tuple y$，共有の受理文脈 $E_0$，任意の試験文脈 $E_1$ を取る．最初の二語 $E_0[\tuple x],E_0[\tuple y]$ から，両方の組の内部の窓はすべて許可されている．$E_1[\tuple x]$ と $E_1[\tuple y]$ では，端のデータ，固定片の窓，またぐ窓が一致するので，骨格への所属が一致する．可換群の値は，両方の文脈で成分の差の同じ和だけ変わり，それは $\mathsf H$ に入る．
\end{proof}

\begin{remark}[現れる窓]
\label{rem:occurring-windows}
定数 $K$ が証明に入るのは，$L$ の語に現れうる辺の種類数を通してだけである（頻度の区間，色の中の不純な窓のベクトル，巡回の長さ）．したがって $K$ は，$L$ の語に現れる $\ell$-窓の数に置き換えられる．これは注意~\ref{rem:degree} で，装飾のあとの次数を改善するのに使う．
% TODO(verification): 証明中の K のすべての出現を確認する（research/README.md，義務2）．
\end{remark}

%======================================================================
\section{装飾}
\label{sec:decoration}

$\mathcal A=(Q,q_0,\delta)$ を $\Sigma$ 上の完全な決定性オートマトンとし，その正の遷移半群を $T$ とする．到達不能な状態は除く．語 $w$ の $p\in Q$ からの\emph{足跡} $D_p(w)$ は，状態 $q$ で $a$ を読むときに $(q,a)$ を出力する．これは $\Gamma=Q\times\Sigma$ 上の語で，文字ごとの射影 $\pi(q,a)=a$ は長さを保つ．$\Gamma$ 上の語は，隣り合う文字が遷移と合っているとき\emph{整合的}であるという．

\begin{definition}[復元]
\label{def:recovery}
$\rho:\Sigma^+\to\mathsf M$ を有限半群への全射準同型とし，$d\ge1$ とする．任意の到達可能な状態 $p$ と任意の正の語 $w$ について，データ
\[
  \bigl(p,\ p\cdot w,\ c_\Gamma(D_p(w)),\ b_d(D_p(w))\bigr)
\]
が $\rho(w)$ を決めるとき，$\mathcal A$ は\emph{幅 $d$ で $\rho$ を復元する}という．
\end{definition}

\begin{mainthm}[装飾の移送；定理~\ref{thm:transfer}]
\label{thm:B}
$\mathcal A$ が幅 $d$ で $\rho$ を復元し，正の語の $L$ への所属が $(\rho(w),N_\ell(w))$ だけで決まり，$s\preceq(c_\Sigma,b_\ell,\rho)$，$L\in\Cmcf{f}{s}$ とする．$\mathcal A$ の遷移半群が局所群なら，$L\cap\Sigma^+$ の任意の reduced good，fan-out $\le f$，rank $\le r$ の表示 $G$ について，$D=f|Q|k^m$，$m=\max(2,d,\ell)$ として $\tau(G,s)=O(v^{2D+3}\bth_G^{D+1})$ が成り立つ．
\end{mainthm}

\begin{theorem}
\label{thm:transfer}
定理~\ref{thm:B} の仮定のうち，$T\in\LG$ の代わりに次だけを仮定する：各非終端 $A$ と，$A$ で実現可能な入口・出口の対のベクトル $(p_i,q_i)_i$ について，各成分が $p_i$ から $q_i$ へ走る $A$ の生成組で，総長が $\beta$ 以下のものがある．このとき
\[
 \tau(G,s)=O\bigl(v^{D'+2}\max(1,\theta,\beta)^{D'+1}\bigr),\qquad
 D'=f(|Q|k)^m,\ \theta=\bth_G .
\]
$T\in\LG$ なら $\beta=O(v\theta)$ であり，指数 $D'$ は $D=f|Q|k^m$ に置き換えられる．これが定理~\ref{thm:B} を与える．
\end{theorem}

\begin{proof}
$L^\#=D_{q_0}(L\cap\Sigma^+)$，$\Gamma$ 上で $t=(c_\Gamma,b_m)$ とおく．

\emph{第1段：$L^\#$ は $m$-局所計数である．}
$\Gamma$ 上の語が $q_0$ からのある語の足跡であるのは，最初の文字の状態が $q_0$ で，整合的であるときに限る．データ $(b_m,N_m)$ は，最初と最後の文字と，より短いすべての因子の個数を決める（長い窓の接頭辞を数え，最後の短い窓を一つ足す．短い語はそのまま読む）．これで整合性を判定でき，$c_\Gamma$，$b_d$，両端の状態が得られ，復元によって $\rho$ が得られる．長さ $\ell$ の個数を射影すれば，射影の $N_\ell$ が得られる．整合的な成分では，$t$ だけからその端の状態，射影の $(c_\Sigma,b_\ell)$，復元によって $\rho$，したがって射影の $s$-型が得られる．

\emph{第2段：$L^\#\in\Cmcf{f}{t}$．}
arity $1\le a\le f$，成分ごとに $t$ が等しい $\Gamma$ 上の正の組 $\tuple X,\tuple Y$，共有の受理文脈 $C_0$，任意の試験文脈 $C_1$ を固定する．四つの語は $C_0[\tuple X],C_0[\tuple Y],C_1[\tuple X],C_1[\tuple Y]$ である．最初の二つから，各成分は整合的で端の状態が一致する．射影した組 $\tuple x=\pi\tuple X$，$\tuple y=\pi\tuple Y$ の $s$-型は等しく，射影した語 $(\pi C_0)[\tuple x],(\pi C_0)[\tuple y],(\pi C_1)[\tuple x],(\pi C_1)[\tuple y]$ には $L$ の代入可能性が使える．$C_1[\tuple X]$ の整合性が固定片の中，先頭，境界で崩れるなら，成分の端の状態が一致するので $\tuple Y$ でも同時に崩れる．そうでなければ，所属は射影の所属である．

\emph{第3段：装飾した文法．}
$G^\#$ の記号を $[A;(p_i,q_i)_i]$ とする．元の各規則と入口・出口のベクトルの各選び方について，各出力テンプレートをその入口状態から走査する．装飾した終端記号を出力し，各変数ではその子の宣言した入口を要求して出口から続け，最後に宣言した出口を確かめる．新しい単項の開始記号で $q_0$ を各終状態につなぐ．導出についての帰納法により，装飾した記号は，宣言した端の状態を持つ $A$ の生成組の足跡をちょうど生成する．trim すると $G^\#$ は reduced good で，fan-out $\le f$，rank $\le r$，記号数 $v'\le1+v|Q|^{2f}$ であり，$L^\#$ を表示する．

\emph{第4段：厚さ．}
装飾した規則が書く終端記号は $\theta$ 個以下で，残った装飾記号はどれも長さ $\beta$ 以下の生成組を持つ．残った規則を目撃するには，$G^\#$ の依存グラフの単純路でその親まで行き，兄弟と子をすべてそうした組で埋めればよい．したがって $\bth_{G^\#}\le(v'+1)(\theta+r\beta)$ である．定理~\ref{thm:L} の証明は，厚さを，規則ごとの終端記号の数と記号ごとの生成組一つを抑えるためにしか使わないので，$\max(1,\theta,\beta)$ を代入してよい．

\emph{第5段：$L$ へ戻す．}
$G$ の実現された型付き非終端 $X$ と，成功する導出の中の生成組 $\tuple x$ を取る．文の足跡から，装飾した記号の生成組 $\tuple X$ と装飾した受理文脈 $F$ が得られる．定理~\ref{thm:L} から，$t$ も装飾した分布も同じで，述べた長さの組 $\tuple Y$ が得られる．$F[\tuple Y]$ は受理されるので，$\tuple Y$ は整合的で，$\pi\tuple Y$ は $\tuple x$ の $s$-型を持つ．任意の試験 $E$ について，四つの語 $(\pi F)[\tuple x],(\pi F)[\pi\tuple Y],E[\tuple x],E[\pi\tuple Y]$ から $\pi\tuple Y\equiv_{L,s}\tuple x$ が従う．装飾が $E$ に延びる必要はない．$\pi$ で長さは変わらない．

$T\in\LG$ なら，実現可能な入口・出口のベクトルは遷移の値を指定した生成組で目撃でき，局所群の短い生成組の定理 \cite[定理 \textit{local-group-polydata}]{KuriyamaCompendium} から長さ $O(v\theta)$ のものが得られる．最後に，$L^\#$ の語に現れるのは整合的な装飾窓だけで，それは $|Q|k^m$ 個以下である．注意~\ref{rem:occurring-windows} により指数は $D=f|Q|k^m$ になる．
\end{proof}

\begin{remark}[この段での局所群の必要性]
\label{rem:lg-needed-for-witnesses}
すべての表示について一様に $\beta$ を抑えるには，$T\in\LG$ が必要である．$T\notin\LG$ なら，$ez=ze=z$ となる $T$ の相異なる冪等元 $e,z$，それを表す正の語 $u,v_0$，$je=j\ne jz$ となる状態 $j$，$q_0$ から $j$ への語 $h$ を選ぶ．fan-out 1，rank 2 の文法 $A_0\to u\mid v_0$，$A_i\to A_{i-1}A_{i-1}\mid v_0$，$S\to SS\mid a\ (a\in\Sigma)\mid hA_i\ (0\le i\le n)$ は reduced good で，大きさ $O(n)$，厚さは有界だが，$j$ から $j$ へ走る $A_n$ の生成組は $u^{2^n}$ だけである．ここでの標的は $\Sigma^+$ で，$\tau$ は定数である．失敗するのは方法で，データではない．
\end{remark}

%======================================================================
\section{復元を許す骨格}
\label{sec:algebra}

\begin{lemma}
\label{lem:lg-closed}
$\LG*\LG=\LG$．
\end{lemma}

\begin{proof}
$U,T\in\LG$ による花輪積 $U\wr_QT$ の冪等元 $E=(F,b)$ と，その局所モノイドの冪等元 $Z=(H,z)$ を取る．$z$ は群 $bTb$ の冪等元なので $z=b$．$ub=u$ となる状態 $u$ では $F(u)$ と $H(u)$ は冪等元で，$EZ=ZE=Z$ から $H(u)$ は群 $F(u)UF(u)$ に入る．よって $H(u)=F(u)$．任意の $u$ では $ub$ が $b$ で固定されるので，$H(u)=F(u)H(ub)=F(u)F(ub)=F(u)$．したがって局所モノイドの冪等元は一つで，群である．
\end{proof}

\begin{mainthm}[復元と $\LECom$；定理~\ref{thm:recovery}，\ref{thm:SLLG-LECom}]
\label{thm:C}
正則な $R\subseteq\Sigma^+$ の統語射を $\rho_R$，統語半群を $S_R$ とする．次は同値である．
\begin{enumerate}[label=(\roman*),leftmargin=*]
\item 遷移半群が $\LG$ に入るあるオートマトンが，ある幅で $\rho_R$ を復元する．
\item $S_R\in\SL*\LG$．
\item $S_R\in\LECom$．
\end{enumerate}
\end{mainthm}

\begin{theorem}
\label{thm:recovery}
定理~\ref{thm:C} の (i) と (ii) は同値である．
\end{theorem}

\begin{proof}
(i)$\Rightarrow$(ii)．初期状態から，足跡の内容は有限の半束で，端の窓は局所自明な半群で累積でき，どちらも装飾した文字を読む．端の状態は下の機械から得られる．したがって $R$ は，$\SL$-over-$\LG$ と $\LI$-over-$\LG$ の連結の直積で認識される．$\LI\subseteq\LG$，補題~\ref{lem:lg-closed}，$\SL*\LG\supseteq\LG$ から $S_R\in\SL*\LG$．

(ii)$\Rightarrow$(i)．$\rho_R$ を文字ごとに，$U\in\SL$，$T\in\LG$ による $U\wr_{T^1}T$ の中の約数表示に持ち上げる．$a$ の持ち上げを $(F_a,t_a)$，$n=|T|$ とする．$T$ の $n$ 個以上の元の積は $p=pt$ となる因子を持つので $T\,E(T)\,T$ に入り，$T$ の冪等元はすべて極小イデアル $K$ に入るので，これは $K$ に含まれる．$p\in K$，$x\in T^1$ なら $xp\le_{\mathcal L}p$ かつ $xp\,\mathcal J\,p$ なので $xp\,\mathcal L\,p$ であり，$hxp=p$ となる $h\in T^1$ がある．

下の状態 $x\mapsto xt_a$ とともに，各事象 $(x,a)$ ごとに，その事象の直後に1増える法 $n+1$ のカウンタを走らせる．遷移半群は $T$ の上の有限可換群の花輪積を割り切り，補題~\ref{lem:lg-closed} から $\LG$ に入る．幅 $n+1$（2以上）を使う．任意の入口から読んだ成分について，最初の $n$ 文字は窓に見える．各事象について，内容の中に現れる増分前の剰余の種類数は，カウンタの初期値によらず $\min(n+1,N_{\rm event})$ である．最初の $n$ 文字の中の出現（$n$ 個以下）を引けば，その事象がそのあとに現れるかが分かる．$x$ を実際の入口，$p$ を最初の $n$ 文字の積，$h$ を $hxp=p$ となるものとする．仮想の入口 $y\in T^1$ については，最初の $n$ 文字を直接模倣する．そのあと，観測された事象 $(z,a)$ での仮想の直前状態は $yhz$ である．間の積 $t$ について $z=xpt$ だからである．$U$ は半束なので，$y$ での上の値は，接頭部の値と，そのあとの相異なる事象についての $F_a(yhz)$ の積である．下の値は $hz_{\rm end}$ である．これで持ち上げのすべての座標，したがって $\rho_R(w)$ が復元される．$n$ より短い語はそのまま読む．
\end{proof}

\begin{theorem}
\label{thm:SLLG-LECom}
$\SL*\LG=\LECom$．
\end{theorem}

\begin{proof}
(a) $\LG=\Gp*\Dd$：自明でない群の擬多様体 $\mathbf H$ について $\mathbf H*\Dd=\mathbf L\mathbf H$ である \cite[例1.3，定理2.8]{TherienWeiss1985}．

(b) $\SL*\Gp=\ECom$．冪等元が可換な有限モノイドはどれも有限逆モノイドを割り切る \cite{Ash1987}．有限逆モノイドはどれも有限の $E$-unitary な被覆を持ち \cite{McAlister1974}，有限の $E$-unitary な逆モノイドはどれも有限の半束と有限群の半直積に埋め込める \cite{OCarroll1976}．よって $\ECom\subseteq\SL*\Gp$．逆に，半束と群の半直積の冪等元は群成分が自明で互いに可換であり，$\ECom$ は擬多様体である \cite{Ash1987}．\cite{MargolisPin1987} も参照．

(c) $\ECom*\Dd=\LECom$．6元の Brandt モノイド $B_2^1$ は逆モノイドなので $\ECom$ に入る．\cite[定理1.5]{TherienWeiss1985} により，グラフ合同 $\beta$ が $\ECom$-合同であることと，集約モノイド $M_\beta=P/\beta\cup\{0,1\}$（つながる路は合成し，ほかの積はすべて $0$）が $\ECom$ に入ることは同値である．$M_\beta$ の冪等元は $0$，$1$，ループの類である．違う頂点のループの積はどちらの順でも $0$ で，同じ頂点のループは，$\beta$ が局所的に $\ECom$-合同なら可換である．したがって局所的に $\ECom$-合同なものはすべて $\ECom$-合同であり，\cite[定理2.8]{TherienWeiss1985} から $\ECom*\Dd=\LECom$．

(a)，結合律 \cite{Tilson1987}，(b)，(c) から
$\SL*\LG=\SL*(\Gp*\Dd)=(\SL*\Gp)*\Dd=\ECom*\Dd=\LECom$．
\end{proof}

\begin{remark}
包含 $\SL*\LG\subseteq\LECom$ は直接にも従う．上が半束の花輪積では，$E=(F,b)$ での局所モノイドの冪等元 $Y,Z$ の下の座標は $b$ で，$b$ で固定される状態では上の座標が可換であり，$EYZ=YZ$，$EZY=ZY$ がこれをすべての状態に伝える．局所的な擬恒等式は約数と積に伝わる．
\end{remark}

\begin{remark}[閾値つきの計数]
\label{rem:LTT}
$\LT=\SL*\Dd$ \cite{BrzozowskiSimon1973,McNaughton1974,Straubing1985} であり，$\LTT$ の統語的な特徴づけは Straubing \cite{Straubing1985} と Th\'erien--Weiss \cite{TherienWeiss1985} の結果を組み合わせたものである（\cite{BeauquierPin1991} 参照）．花輪積の形では $\LTT=\ACom*\Dd$．$\SL\subseteq\LT\subseteq\LTT$ と $\ACom\subseteq\ECom=\SL*\Gp$ から
\[
  \SL*\LG=\LT*\LG=\LTT*\LG .
\]
包含 $\LTT*\Gp\subsetneq\SL*\LG$ は真である．因子 $\mathtt{ab}$ の個数が偶数の語の言語の統語半群は $\LG$ に入る（直前の文字を右零半群で保ち，$\mathtt{ab}$ で群のビットを反転する）．しかし，非周期的な層を群の上に置くどんな認識でも，$N$ を下の指数の倍数として $(\mathtt a^N\mathtt b^N)^j$ の持ち上げの上の座標はやがて一定になるのに，所属は $j$ とともに交互に変わる．
\end{remark}

%======================================================================
\section{主定理と系}
\label{sec:main}

\begin{theorem}
\label{thm:main}
定理~\ref{thm:A} が成り立つ．より正確には，$S_R\in\LECom$ なら，遷移半群が $\LG$ に入るオートマトン $\mathcal A$ と幅 $d$ があって $\mathcal A$ は幅 $d$ で $\rho_R$ を復元し，定理~\ref{thm:A} の仮定のもとで
\[
  \tau(G,s)=O\bigl(v^{2D+3}\bth_G^{D+1}\bigr),\qquad D=f|Q|k^{\max(2,d,\ell)}
\]
が成り立つ．
\end{theorem}

\begin{proof}
定理~\ref{thm:SLLG-LECom} から $S_R\in\SL*\LG$ なので，定理~\ref{thm:recovery} が $\mathcal A$ と $d$ を与える．$\rho=\rho_R$ として定理~\ref{thm:B} を，特性データについては事実~\ref{fact:char-data} を適用する．
\end{proof}

次数は，$S_R$ の花輪積分解から作った復元オートマトンに依存する．次数を最小化することは試みない．

\begin{corollary}[群の値]
\label{cor:group}
$\varphi:\Sigma^+\to\mathsf G$ を有限群への全射準同型とする．所属が $(b_\ell(w),N_\ell(w),\varphi(w))$ だけで決まり，$s\preceq(c_\Sigma,b_\ell,\varphi)$，$L\in\Cmcf{f}{s}$ なら，$D=f|\mathsf G|k^{\max(2,\ell)}$ として $\tau(G,s)=O(v^{2D+3}\bth_G^{D+1})$．
\end{corollary}

\begin{proof}
$\mathsf G$ の右正則作用は幅 1 で $\varphi$ を復元する．入口 $p$ から出口 $q$ へ読んだ成分の値は $p^{-1}q$ である．その遷移半群は群である．定理~\ref{thm:B} を適用する．射影の窓の型と因子の個数は $b_m$ と $N_m$ から読む．
\end{proof}

\begin{remark}[次数]
\label{rem:degree}
定理~\ref{thm:L} の色に群の値と辺の個数の $|\mathsf G|$ を法とした値を加え，大きな多重集合の積についての Dénes--Hermann 型の補題 \cite{DenesHermann1982} を使う直接の議論は，系~\ref{cor:group} でよりよい指数 $fk^\ell$ を与える（研究記録 R14，\texttt{research/2026-11-23\_R14\_astra\_automata\_skeleton}）．装飾は整合性を確かめるのに幅2以上を要するので，この指数は定理~\ref{thm:B} 経由では回復しない．
% TODO(paper): 直接の証明を付録として入れるか決める．
\end{remark}

\begin{example}[偶数 run]
\label{ex:even-runs}
$E\subseteq\{\mathtt a,\mathtt b\}^+$ を，極大な $\mathtt a$-run の長さがすべて偶数である語の集合とする．その最小オートマトンは二つの生きた状態（今の run の偶奇）と死状態を持ち，$\mathtt a$ は生きた状態を入れ替え，$\mathtt b$ は $0$ を固定して $1$ を殺す．正の統語半群は，二元集合の部分単射七個からなる対称逆モノイド $I_2$ で，代表元は $\mathtt{aa},\mathtt a,\mathtt b,\mathtt{ba},\mathtt{ab},\mathtt{aba},\mathtt{bab}$ である．$\mathtt bI_2\mathtt b=\{\mathtt b,0\}\cong U_1$ なので $I_2\notin\LG$，$I_2$ は逆モノイドなので $I_2\in\ECom\subseteq\LECom$．偶奇のオートマトン（位数2の群）は幅 1 で $\rho_E$ を復元する．入口と出口の偶奇と，装飾された $\mathtt b$ が読まれる偶奇の集合から，成分の部分写像が得られる．したがって定理~\ref{thm:A} は，たとえば $E\cap g^{-1}(\gamma+\mathsf H)$ が $f$-MCFL ならそれに適用できる．像が $\LG$ に入るどの型付けでも，$E$ は1次元の代入可能性さえ満たさない．$e=s(\mathtt a^{2h})$ が冪等元で，$d$ を群 $e\mathsf Se$ の中の $e\,s(\mathtt b)\,e$ の位数とすると，組 $\mathtt a^{2h}$ と $(\mathtt a^{2h}\mathtt b\mathtt a^{2h})^d$ は型 $e$ を持ち空の文脈を共有するが，$\mathtt a\blank\mathtt a$ は前者だけを受理する．
\end{example}

\begin{corollary}[局所計数と路の標的]
\label{cor:local}
定理~\ref{thm:L} は一状態のオートマトンの場合である．特に，有向グラフについて，文字を辺として読んだとき両端の頂点を固定した路になる語の言語に，文字の個数についての任意の条件を交わらせたものでは，$\tau$ は全 fan-out で多項式である（$\ell=2$ の場合）．
\end{corollary}

%======================================================================
\section{方法の限界}
\label{sec:limits}

\begin{theorem}
\label{thm:limits}
$S_R\notin\LECom$ とする．遷移半群が $\LG$ に入る任意のオートマトン $\mathcal A$，任意の幅 $m$，$q_0$ からの足跡の上の $(b_m,N_m)$ の値の任意の集合 $\mathcal Q$ について，集合 $\{w:(b_m,N_m)(D_{q_0}(w))\in\mathcal Q\}$ は $R$ と異なる．
\end{theorem}

\begin{proof}
$\mathcal A$ の上に，最後の $m-1$ 個の装飾文字（または短い接頭辞）を保つレジスタを積む．レジスタの正の像は局所自明なので，新しい下の機械 $\mathcal A'$ の遷移半群 $T'$ もまた $\LG$ に入り（補題~\ref{lem:lg-closed}），各段で新しく完成した装飾窓を出力できる．最終の端の窓はレジスタにある．

$S_R\notin\LECom$ なので，$\rho_R(c_0)$ が冪等元で，$\rho_R(c_0uc_0)$ と $\rho_R(c_0vc_0)$ が冪等元で可換でない語 $c_0,u,v$ がある．$u,v$ を $c_0uc_0,c_0vc_0$ に替えて，$u,v$ は $e=\rho_R(c_0)$ の局所モノイドの冪等元で $\rho_R(uv)\ne\rho_R(vu)$ としてよい．$\alpha c_0uvc_0\beta$ と $\alpha c_0vuc_0\beta$ を分ける文脈 $\alpha\blank\beta$ を選ぶ．$C=c_0^N$ が $T'$ の冪等元 $e'$ として働くように $N$ を取る．$CuC$ と $CvC$ は群 $e'T'e'$ の中で働くので，共通の正の冪 $j$ について $U=(CuC)^j$ と $V=(CvC)^j$ はともに $e'$ として働く．$\alpha C$ を読んだあと，$U$ と $V$ は $\mathcal A'$ の同じ状態でのループなので，$\alpha CUV\beta$ と $\alpha CVU\beta$ は最終の窓も，完成した窓の個数も等しい（各ループは決まったベクトルを足す）．統語値はそれぞれ同じ文脈の中に $uv$ と $vu$ を含むので，ちょうど一方だけが $R$ に入る（$u,v$ は局所モノイドの冪等元なので，冪は値を変えない）．
\end{proof}

\begin{example}[$\LECom$ の外の最小の骨格]
\label{ex:minimal}
\begin{enumerate}[label=(\alph*),leftmargin=*]
\item 最後の $\mathtt a$-run の長さが偶数：$\{\mathtt a,\mathtt b\}$ 上の $(\mathtt{aa})^+\cup\Sigma^*\mathtt b(\mathtt{aa})^*$．変換モノイドは全変換モノイド $T_2$ で，$\mathtt{aa}=1$ での局所モノイドの冪等元 $\mathtt b$ と $\mathtt{ba}$ は相異なる定値写像である．
\item $\mathtt c$ を中立とする $\{\mathtt a,\mathtt b,\mathtt c\}$ 上の散在パターン $\mathtt a\cdots\mathtt b$：$\mathtt c$ は統語的な単位元，$\mathtt a$ と $\mathtt b$ は冪等で，$\mathtt{ab}$ は受理，$\mathtt{ba}$ は受理されない．（$\{\mathtt a,\mathtt b\}$ だけなら同じパターンは局所テスト可能である．）
\item 恒等の文字を持つ flip-flop：$r_0r_1=r_1\ne r_0=r_1r_0$ のモノイド $\{1,r_0,r_1\}$．位数2以下の半群はどれも $\LG$ か $\SL$ に入るので，$\LECom$ の外の最小の位数は3である．
\end{enumerate}
\end{example}

\begin{remark}[射影]
\label{rem:projections}
例~\ref{ex:minimal}(b) は，先行の研究記録 \cite{KuriyamaCompendium} の別の議論（依存する文字の非交互性の定理と，独立な文字ブロックの定理）で扱われる．そこでは部分アルファベットへの消去射影を使う．消去射影は，長さを保つ射影を使う定理~\ref{thm:B} の対象外である．両者を合わせた主張は未解決である．
\end{remark}

\begin{remark}[限界が言わないこと]
定理~\ref{thm:limits} と注意~\ref{rem:lg-needed-for-witnesses} は装飾の方法の限界であって，$\tau$ の下界ではない．骨格 $R$ だけを型付け $\rho_R$ で扱うなら，$\tau$ は定数である．超多項式の下界には，正則でない計数条件と，正しい分布を持つ組がどれも長くなる表示の族が要る．
\end{remark}

\begin{openproblem}[{\cite[未解決問題 \textit{fanout-two}]{KuriyamaCompendium}}]
\label{op:fanout-two}
$f\ge2$ について，$\LG$ に含まれないどの擬多様体も，クラス $\Cmcf{f}{s}$ が総長多項式の特性データを持たないような型付けを含むか．
\end{openproblem}

本論文の結果から，肯定の答えには，骨格が正則でないか統語半群が $\LECom$ の外にある標的，または計数条件が有界長の因子についてのものでない標的を使わなければならない．

%======================================================================
\section{結論}
\label{sec:conclusion}

型付けについての局所群の境界に対応する，標的の側の結果を与えた．$\LECom$ に入る正則な骨格と，有界長の因子についての任意の計数条件で作られる標的は，$\LG$ の外の型付けについても，全 fan-out で多項式の代表長を持つ．証明は，局所群の装飾によって標的を局所計数言語に移し，その装飾はちょうど $\LECom=\SL*\LG$ について存在する．最適な次数，消去射影との統合（注意~\ref{rem:projections}），計数条件を伴う $\LECom$ の外の骨格，未解決問題~\ref{op:fanout-two} が残る問題である．

\paragraph{謝辞．}
% TODO(paper): 雑誌の方針に従い，AI 利用の申告と謝辞を書く．

%======================================================================
\begin{thebibliography}{99}

\bibitem{Ash1987}
C.~J.~Ash.
\newblock Finite semigroups with commuting idempotents.
\newblock {\em Journal of the Australian Mathematical Society (Series A)},
43:81--90, 1987.

\bibitem{BeauquierPin1991}
D.~Beauquier and J.-E.~Pin.
\newblock Languages and scanners.
\newblock {\em Theoretical Computer Science}, 84(1):3--21, 1991.

\bibitem{BrzozowskiSimon1973}
J.~A.~Brzozowski and I.~Simon.
\newblock Characterizations of locally testable events.
\newblock {\em Discrete Mathematics}, 4:243--271, 1973.

\bibitem{CostaNogueiraTeixeira2021}
J.~C.~Costa, C.~Nogueira, and M.~L.~Teixeira.
\newblock The word problem for $\kappa$-terms over the pseudovariety of local
groups.
\newblock {\em Semigroup Forum}, 103:439--468, 2021.

\bibitem{DenesHermann1982}
J.~D\'enes and P.~Hermann.
\newblock On the product of all elements in a finite group.
\newblock {\em Annals of Discrete Mathematics}, 15:105--109, 1982.

\bibitem{KuriyamaTCS1}
T.~Kuriyama.
\newblock Distributional learning of context-free languages under fixed
finite-semigroup typing.
\newblock arXiv:1409.6247v5, 2026.

\bibitem{KuriyamaMCFGpreprint}
T.~Kuriyama.
\newblock The value of finite observation in positive-data learning of multiple
context-free languages.
\newblock arXiv:2605.11644v3, 2026.

\bibitem{KuriyamaCompendium}
T.~Kuriyama.
\newblock Multiple context-free languages under fixed finite-semigroup typing:
research compendium.
\newblock \url{https://github.com/growupkuriyama-hub/Papers/tree/main/02-v4-fixed-s-mcfg},
2026.

\bibitem{KuriyamaTypingClass}
T.~Kuriyama.
\newblock Which fixed typings give polynomial characteristic data? A
classification for substitutable context-free languages.
\newblock Manuscript, 2026.

\bibitem{MargolisPin1987}
S.~W.~Margolis and J.-E.~Pin.
\newblock Inverse semigroups and extensions of groups by semilattices.
\newblock {\em Journal of Algebra}, 110:277--297, 1987.

\bibitem{McAlister1974}
D.~B.~McAlister.
\newblock Groups, semilattices and inverse semigroups.
\newblock {\em Transactions of the American Mathematical Society}, 192:227--244,
1974.

\bibitem{McNaughton1974}
R.~McNaughton.
\newblock Algebraic decision procedures for local testability.
\newblock {\em Mathematical Systems Theory}, 8:60--76, 1974.

\bibitem{OCarroll1976}
L.~O'Carroll.
\newblock Embedding theorems for proper inverse semigroups.
\newblock {\em Journal of Algebra}, 42:26--40, 1976.

\bibitem{SekiEtAl1991}
H.~Seki, T.~Matsumura, M.~Fujii, and T.~Kasami.
\newblock On multiple context-free grammars.
\newblock {\em Theoretical Computer Science}, 88(2):191--229, 1991.

\bibitem{Straubing1985}
H.~Straubing.
\newblock Finite semigroup varieties of the form $\mathbf V*\mathbf D$.
\newblock {\em Journal of Pure and Applied Algebra}, 36:53--94, 1985.

\bibitem{TherienWeiss1985}
D.~Th\'erien and A.~Weiss.
\newblock Graph congruences and wreath products.
\newblock {\em Journal of Pure and Applied Algebra}, 36:205--215, 1985.

\bibitem{Tilson1987}
B.~Tilson.
\newblock Categories as algebra: an essential ingredient in the theory of
monoids.
\newblock {\em Journal of Pure and Applied Algebra}, 48:83--198, 1987.

\bibitem{Yoshinaka2011}
R.~Yoshinaka.
\newblock Efficient learning of multiple context-free languages with
multidimensional substitutability from positive data.
\newblock {\em Theoretical Computer Science}, 412(19):1821--1831, 2011.

\end{thebibliography}

\end{document}
